% ─────────────────────────────────────────────────────────────────────────────
%  Phụ lục
%  Hướng dẫn của Khoa: phân chia công việc nhóm (nếu làm nhóm), tài liệu tham
%  khảo và trang web lấy mã nguồn, các ký hiệu viết tắt.
%  Danh mục viết tắt đã đặt ở đầu quyển, tài liệu tham khảo in trước phụ lục.
%
%  Chương "Phân chia công việc" đã bỏ vì đây là báo cáo cá nhân — chương đó chỉ
%  bắt buộc với báo cáo làm theo nhóm.
% ─────────────────────────────────────────────────────────────────────────────
\chapter{Hướng dẫn cài đặt và chạy}
\label{app:caidat}

\section{Yêu cầu môi trường}

Python 3.10 trở lên, Node.js 18 trở lên. Không bắt buộc có GPU; hệ thống tự chọn
CUDA nếu máy có card NVIDIA, MPS trên chip Apple, còn lại chạy CPU.

\section{Khởi động máy chủ}

\begin{lstlisting}[language=bash, caption={Cài đặt và chạy backend}]
cd backend
python -m venv .venv
source .venv/bin/activate        # Windows: .venv\Scripts\activate
pip install -r requirements.txt
python run.py
\end{lstlisting}

Máy chủ lắng nghe tại \code{http://localhost:8000}, tài liệu API tại
\code{http://localhost:8000/docs}. Lần chạy đầu tiên, thư viện Ultralytics tự tải
tệp trọng số \code{yolov8n.pt} khoảng 6~MB.

\section{Khởi động giao diện}

\begin{lstlisting}[language=bash, caption={Cài đặt và chạy frontend}]
cd frontend
npm install
npm run dev
\end{lstlisting}

Mở trình duyệt tại \code{http://localhost:3000}. Vite chuyển tiếp mọi lời gọi
\code{/api} sang máy chủ ở cổng 8000, nên không cần cấu hình CORS khi phát triển.

\section{Thêm camera đầu tiên}

Vào trang \emph{Quản trị}, thẻ \emph{Nguồn camera}. Nhập tên, chọn loại nguồn rồi
điền đường dẫn: số \code{0} cho webcam tích hợp, chuỗi \code{rtsp://\dots} cho
camera IP, hoặc đường dẫn tuyệt đối tới tệp video. Bấm \emph{Kiểm tra kết nối}
trước khi lưu.

Sang trang \emph{AI Pipeline}, chọn camera vừa thêm, chọn bài toán \emph{Đếm đối
tượng ra/vào}, kéo vạch trên ảnh, rồi bấm \emph{Lưu và chạy ngay}.

\chapter{Lược đồ cơ sở dữ liệu}
\label{app:schema}

\begin{table}[H]
\centering
\caption{Các bảng chính trong cơ sở dữ liệu SQLite.}
\label{tab:bangcsdl}
\small
\begin{tabularx}{\textwidth}{@{}lX@{}}
\toprule
\textbf{Bảng} & \textbf{Nội dung lưu trữ} \\
\midrule
\code{cameras}   & Danh sách camera: tên, vị trí, nguồn video, trạng thái bật/tắt \\
\code{pipelines} & Cấu hình bài toán: loại bài toán, lớp đối tượng, toạ độ vạch \\
\code{events}    & Sự kiện vượt vạch: thời điểm, chiều, mã vết, ảnh chụp, trạng thái xử lý \\
\code{counts}    & Mốc số đếm ghi định kỳ, dùng dựng biểu đồ lưu lượng \\
\code{segments}  & Chỉ mục các đoạn video đã ghi: thời điểm bắt đầu, thời lượng, đường dẫn \\
\code{logs}      & Nhật ký hoạt động của hệ thống \\
\code{settings}  & Cấu hình dạng khoá--giá trị, ví dụ thời hạn lưu trữ \\
\bottomrule
\end{tabularx}
\end{table}

\chapter{Danh sách điểm cuối API}
\label{app:api}

Máy chủ cung cấp 40 điểm cuối chia theo bảy nhóm. Tài liệu chi tiết --- kiểu dữ liệu
vào, kiểu trả về, ví dụ phản hồi --- nằm ở tệp \code{docs/api.md} trong mã nguồn, được
\emph{sinh tự động} từ hệ thống đang chạy bằng \code{scripts/sinh\_tai\_lieu\_api.py}.
Sinh tự động thay vì viết tay để tài liệu không lệch khỏi mã nguồn --- một lỗi dự án
này đã mắc ba lần với các tham số khác.

Đặc tả OpenAPI đầy đủ có sẵn tại \code{/openapi.json}, giao diện thử trực tiếp tại
\code{/docs}.

\begin{table}[H]
\centering
\caption{Các nhóm điểm cuối API.}
\label{tab:nhomapi}
\small
\begin{tabularx}{\textwidth}{@{}lcX@{}}
\toprule
\textbf{Nhóm} & \textbf{Số lượng} & \textbf{Chức năng} \\
\midrule
\code{/api/cameras}  & 10 & Khai báo nguồn video, luồng MJPEG, ảnh chụp, trạng thái thực \\
\code{/api/pipelines} &  9 & Cấu hình bài toán AI, chạy và dừng, đặt lại số đếm \\
\code{/api/events}   &  3 & Sự kiện vượt vạch, xác nhận xử lý, ảnh chụp sự kiện \\
\code{/api/search}   &  1 & Tìm kiếm bằng câu tiếng Việt \\
\code{/api/playback} &  6 & Dòng thời gian, đoạn ghi hình, phát và tải về \\
\code{/api/admin}    &  7 & Cấu hình hệ thống, dung lượng, dọn dữ liệu, nhật ký \\
\code{/api/ai}       &  2 & Bộ mô tả công cụ và cổng gọi cho trợ lý AI \\
\bottomrule
\end{tabularx}
\end{table}

Nhóm \code{/api/ai} là cổng để trợ lý AI bên ngoài gọi vào hệ thống:
\code{GET /api/ai/tools} trả về bộ mô tả công cụ dạng máy đọc được, còn
\code{POST /api/ai/call} thực thi một lời gọi. Bộ công cụ \textbf{chỉ gồm thao tác
đọc} --- trợ lý không thêm, sửa hay xoá được gì, vì mô hình ngôn ngữ hiểu sai một câu
hỏi thì chỉ ra câu trả lời sai, còn hiểu sai một lệnh xoá thì mất dữ liệu.
